\documentclass[12pt]{article}

\usepackage{amsmath, amssymb}
\usepackage{authblk}
\usepackage{geometry}
\usepackage{hyperref}
\geometry{margin=1in}

\title{\textbf{A Quantized Rupture Cosmology:\\
Cyclic Dynamics, Multiverse Branching, and Entropy Inheritance}}

\author[1]{Paul Jarvis}
\affil[1]{Independent Researcher, United Kingdom}

\date{27 May 2026}

\begin{document}
\maketitle

\begin{abstract}
I propose a rupture-driven cyclic cosmology in which black hole formation increases a singularity density variable $\rho_s(t)$ until a critical threshold $\rho_{\rm crit}$ is reached. Crossing this threshold triggers a rupture that ends the cycle and produces $N = 6m$ descendant universes, where $m$ is a quantized excitation mode determined by rupture strength. Numerical simulations show that for natural parameter choices the fundamental mode $m=1$ overwhelmingly dominates, yielding six universes per rupture. Extending the model to a branching universe tree and adding entropy tracking reveals a thermodynamic arrow across generations: each universe produces entropy, dumps it into the singularity at rupture, and passes it forward to its descendants. The resulting structure is a stable, self-similar multiverse with quantized reproduction and inherited entropy.
\end{abstract}

\section{Introduction}

Cyclic cosmologies have been explored in several contexts, including conformal cyclic cosmology (CCC) \cite{Penrose2010}, ekpyrotic and brane-based cyclic models \cite{SteinhardtTurok2002,SteinhardtTurok2005}, and loop quantum cosmology (LQC) bounce scenarios \cite{Ashtekar2006}. A persistent challenge is identifying a physically motivated mechanism that both terminates each cycle and seeds the next.

In this work, I introduce a rupture-driven cyclic cosmology in which black hole formation drives the accumulation of a singularity density variable $\rho_s(t)$. When $\rho_s(t)$ exceeds a critical threshold $\rho_{\rm crit}$, the universe undergoes a rupture that resets the cycle and produces a discrete number of descendant universes. A key feature of the model is that rupture outcomes are quantized: the number of universes produced is always a multiple of six, reflecting an assumed six-directional symmetry of the confining structure.

\section{Relation to Previous Work}

Smolin's cosmological natural selection \cite{Smolin1992,Smolin1997} proposes that black holes generate new universes with varied parameters, forming a branching multiverse. However, Smolin's model is not cyclic, does not include a rupture threshold, and does not impose quantized reproduction.

Penrose's CCC \cite{Penrose2010} describes an infinite sequence of aeons related by conformal rescaling, with entropy effectively reset between cycles. CCC is cyclic but non-branching: each aeon leads to exactly one successor.

Ekpyrotic models \cite{SteinhardtTurok2002,SteinhardtTurok2005} involve brane collisions that periodically reheat the universe, but do not incorporate discrete reproduction or black-hole-driven rupture.

Eternal inflation \cite{Guth2007} produces a branching multiverse but lacks cyclic dynamics, a global rupture threshold, or quantized offspring counts.

Loop quantum cosmology \cite{Ashtekar2006} replaces singularities with bounces but does not introduce branching or reproduction.

To the author's knowledge, no previous model combines: (i) a singularity density accumulator driven by black hole formation, (ii) a critical rupture threshold, (iii) quantized reproduction in multiples of six, (iv) a branching universe tree, and (v) explicit entropy inheritance across generations.

\section{Defining Equation of the Rupture Cosmology}

The essential physics of the model can be expressed in a single compact equation:

\begin{equation}
\boxed{
N(t) = 6\,\Theta\!\left(\rho_s(t) - \rho_{\rm crit}\right)\,
m\!\left(\Delta(t)\right)
}
\label{eq:master}
\end{equation}

where:
\begin{itemize}
\item $N(t)$ is the number of universes produced at rupture,
\item $\Theta$ is the Heaviside step function,
\item $\rho_s(t)$ is the singularity density,
\item $\rho_{\rm crit}$ is the rupture threshold,
\item $\Delta(t) = \rho_s(t) - \rho_{\rm crit}$ is the overshoot,
\item $m(\Delta)$ is a quantized excitation mode.
\end{itemize}

The excitation mode is defined by discrete overshoot intervals:
\begin{equation}
m(\Delta)=
\begin{cases}
1, & \Delta < 500,\

\[4pt]
2, & 500 \le \Delta < 2000,\

\[4pt]
3, & 2000 \le \Delta < 5000,\

\[4pt]
4, & \Delta \ge 5000.
\end{cases}
\label{eq:modes}
\end{equation}

The singularity density evolves according to
\begin{equation}
\boxed{
\frac{d\rho_s}{dt} = \alpha\,\dot{M}_{\rm BH}(t)
}
\label{eq:rho_evolution}
\end{equation}
with black hole formation rate
\begin{equation}
\boxed{
\dot{M}_{\rm BH}(t) = A\, t^p e^{-t/\tau}
}
\label{eq:bh_rate}
\end{equation}

Equations \eqref{eq:master}--\eqref{eq:bh_rate} define the complete rupture cosmology: black hole formation drives the growth of $\rho_s(t)$ until the threshold $\rho_{\rm crit}$ is crossed, at which point the universe ruptures and produces $N = 6m$ descendant universes determined by the quantized excitation mode $m$.

\section{Universe Tree}

Each universe is represented as a node with:
\begin{itemize}
\item generation index,
\item rupture time,
\item entropy produced,
\item entropy inherited,
\item number of children.
\end{itemize}

A multiverse tree is built by recursively evolving each universe and spawning $N=6m$ descendants.

\section{Entropy}

Entropy production is proportional to black hole formation:
\begin{equation}
S_{\rm produced} = \int \dot{M}_{\rm BH}(t)\, dt.
\end{equation}

At rupture:
\begin{equation}
S_{\rm dumped} = S_{\rm produced} + S_{\rm inherited}.
\end{equation}

Each child inherits:
\begin{equation}
S_{\rm inherited}^{\rm child} = S_{\rm dumped}^{\rm parent}.
\end{equation}

This produces a thermodynamic arrow across generations.

\section{Results}

Simulations show:
\begin{itemize}
\item stable rupture times ($t_{\rm rupt} \approx 5.6$ Gyr),
\item dominance of the fundamental mode ($m=1$),
\item linear entropy inheritance,
\item a stable, self-similar multiverse tree.
\end{itemize}

\section{Conclusion}

This rupture-driven cosmology provides a simple mechanism for cyclic evolution, quantized reproduction, and entropy inheritance. The model naturally favors six-universe ruptures, producing a stable branching multiverse with a thermodynamic arrow.

\begin{thebibliography}{9}

\bibitem{Smolin1992}
L.~Smolin,
``Did the Universe Evolve?,'' 
\textit{Class. Quant. Grav.} \textbf{9}, 173--191 (1992).

\bibitem{Smolin1997}
L.~Smolin,
\textit{The Life of the Cosmos},
Oxford University Press (1997).

\bibitem{Penrose2010}
R.~Penrose,
\textit{Cycles of Time: An Extraordinary New View of the Universe},
Bodley Head (2010).

\bibitem{SteinhardtTurok2002}
P.~J.~Steinhardt and N.~Turok,
``A Cyclic Model of the Universe,''
\textit{Science} \textbf{296}, 1436--1439 (2002).

\bibitem{SteinhardtTurok2005}
P.~J.~Steinhardt and N.~Turok,
``The Cyclic Model Simplified,''
\textit{New Astron. Rev.} \textbf{49}, 43--57 (2005).

\bibitem{Guth2007}
A.~H.~Guth,
``Eternal inflation and its implications,''
\textit{J. Phys. A} \textbf{40}, 6811--6826 (2007).

\bibitem{Ashtekar2006}
A.~Ashtekar, T.~Pawlowski and P.~Singh,
``Quantum Nature of the Big Bang,''
\textit{Phys. Rev. Lett.} \textbf{96}, 141301 (2006).

\end{thebibliography}

\end{document}
